\documentclass[11pt,a4paper]{article}

\usepackage[utf8]{inputenc}
\usepackage[T1]{fontenc}
\usepackage[spanish,es-nodecimaldot]{babel}
\usepackage[a4paper,margin=2.5cm,headheight=15pt]{geometry}
\usepackage{amsmath,amssymb}
\usepackage{array}
\usepackage{booktabs}
\usepackage{longtable}
\usepackage{tabularx}
\usepackage{graphicx}
\usepackage{xcolor}
\usepackage{fancyhdr}
\usepackage{caption}
\usepackage{float}
\usepackage{csquotes}
\usepackage{enumitem}
\usepackage{textcomp}
\usepackage{hyperref}

\definecolor{unapblue}{HTML}{1F5D93}
\definecolor{brandblue}{HTML}{578BC5}
\definecolor{brandorange}{HTML}{DB8963}
\definecolor{ink}{HTML}{263445}
\definecolor{softblue}{HTML}{E7F1FE}
\definecolor{softorange}{HTML}{FEECE4}
\definecolor{linegray}{HTML}{DCE3EC}

\hypersetup{
  colorlinks=true,
  linkcolor=unapblue,
  citecolor=unapblue,
  urlcolor=unapblue,
  pdftitle={Informe técnico: Riesgo agroclimático por heladas en la región Puno},
  pdfauthor={Aracayo Mamani, Jhon Marco; Canaza Paucara, Juan Diego; Luque Pacheco, Angie Tatiana},
  pdfsubject={Proyecto final de Visualización de Datos, SIS328},
  pdfkeywords={heladas, Puno, ERA5-Land, dashboard, visualización de datos, PCA}
}
\urlstyle{same}

\pagestyle{fancy}
\fancyhf{}
\fancyhead[L]{\small UNAP -- Visualización de Datos}
\fancyhead[R]{\small E.P. Ingeniería de Sistemas}
\fancyfoot[C]{\thepage}
\renewcommand{\headrulewidth}{0.4pt}
\renewcommand{\footrulewidth}{0pt}

\setlength{\parindent}{1.15em}
\setlength{\parskip}{0.25em}
\raggedbottom
\setlist[itemize]{leftmargin=*,itemsep=0.25em,topsep=0.35em}
\setlist[enumerate]{leftmargin=*,itemsep=0.3em,topsep=0.35em}
\captionsetup{font=small,labelfont=bf,labelsep=period}
\graphicspath{{figuras/}{../assets/capturas/}{../assets/logos/}}

% Archivo generado. No editar a mano.
\newcommand{\NRegistros}{47\,489}
\newcommand{\NVariables}{52}
\newcommand{\NProvincias}{13}
\newcommand{\NDias}{3\,653}
\newcommand{\NHeladas}{7\,620}
\newcommand{\NHeladasAgronomicas}{18\,877}
\newcommand{\TasaHelada}{16,05\,\%}
\newcommand{\TasaHeladaAgronomica}{39,75\,\%}
\newcommand{\FactorAgronomico}{2,48}
\newcommand{\TminMedia}{3,79~\textdegree C}
\newcommand{\TminAbsoluta}{-9,1~\textdegree C}
\newcommand{\RiesgoMedio}{31,71}
\newcommand{\TasaJulio}{49,93\,\%}
\newcommand{\OLSPendiente}{-1,37}
\newcommand{\OLSRdos}{0,260}
\newcommand{\OLSp}{0,132}
\newcommand{\MKTau}{-0,422}
\newcommand{\MKp}{0,107}
\newcommand{\SenPendiente}{-1,00}
\newcommand{\AltitudR}{0,485}
\newcommand{\AltitudP}{0,093}
\newcommand{\LagoTasa}{11,40\,\%}
\newcommand{\RestoTasa}{20,03\,\%}
\newcommand{\FuenteR}{0,611}
\newcommand{\FuenteSesgo}{1,29~\textdegree C}
\newcommand{\FuenteMAE}{2,54~\textdegree C}
\newcommand{\FuenteConcordancia}{83,06\,\%}
\newcommand{\PCAUno}{48,20\,\%}
\newcommand{\PCADosAcum}{68,66\,\%}
\newcommand{\PCANoventa}{5}
\newcommand{\KOptimo}{3}
\newcommand{\Silueta}{0,306}
\newcommand{\NAnomalias}{475}
\newcommand{\RFAUC}{0,990}
\newcommand{\RFAP}{0,948}
\newcommand{\RFExactitud}{95,23\,\%}
\newcommand{\RFSensibilidad}{94,47\,\%}
\newcommand{\RFPrecision}{78,16\,\%}


\newcommand{\projectfigure}[4][0.88\textwidth]{%
  \begin{figure}[H]
    \centering
    \IfFileExists{#2}{%
      \includegraphics[width=#1,height=0.70\textheight,keepaspectratio]{#2}%
    }{%
      \fbox{\parbox[c][4.5cm][c]{0.82\textwidth}{\centering
        Evidencia visual generada por el flujo reproducible del proyecto.\\
        Artefacto esperado: \texttt{\detokenize{#2}}}}%
    }
    \caption{#3}
    \label{#4}
  \end{figure}
}

\newcommand{\smallcapslabel}[1]{\textcolor{unapblue}{\textbf{#1}}}

\begin{document}
\hypersetup{pageanchor=false}

% -----------------------------------------------------------------------------
% Portada institucional
% -----------------------------------------------------------------------------
\begin{titlepage}
  \thispagestyle{empty}
  \begin{center}
    \begin{minipage}{0.19\textwidth}
      \raggedright
      \IfFileExists{../assets/logos/logo_unap.png}{\includegraphics[height=2.3cm]{../assets/logos/logo_unap.png}}{}
    \end{minipage}
    \begin{minipage}{0.58\textwidth}
      \centering
      {\large\bfseries\color{unapblue} UNIVERSIDAD NACIONAL DEL ALTIPLANO}\par
      \vspace{0.25em}
      {\small Escuela Profesional de Ingeniería de Sistemas}\par
      {\small Facultad de Ingeniería Mecánica Eléctrica, Electrónica y Sistemas}
    \end{minipage}
    \begin{minipage}{0.19\textwidth}
      \raggedleft
      \IfFileExists{../assets/logos/logo_episi.png}{\includegraphics[height=2.3cm]{../assets/logos/logo_episi.png}}{}
    \end{minipage}

    \vspace{0.7cm}
    \textcolor{brandorange}{\rule{0.92\textwidth}{1.4pt}}
    \vfill

    {\LARGE\bfseries\color{unapblue}
      RIESGO AGROCLIMÁTICO POR HELADAS\\[0.25em]
      EN LA REGIÓN PUNO\par}

    \vspace{0.8cm}
    {\large\bfseries Dashboard interactivo para el análisis y la visualización de datos orientado a la toma de decisiones\par}

    \vspace{1.0cm}
    {\large Informe técnico presentado para la asignatura de Visualización de Datos\par}

    \vspace{1.0cm}
    \renewcommand{\arraystretch}{1.35}
    \begin{tabular}{>{\bfseries}p{3.2cm}p{9.7cm}}
      Docente: & Ing. Edwin Edgar Mestas Yucra \\
      Asignatura: & Visualización de Datos \\
      Código: & SIS328 \\
      Semestre: & 2026-I \\
      Ciclo: & X \\
      Integrantes: & Aracayo Mamani, Jhon Marco \\
                   & Canaza Paucara, Juan Diego \\
                   & Luque Pacheco, Angie Tatiana \\
      Fuente: & ERA5-Land / ERA5 (ECMWF--Copernicus), con validación MERRA-2 (NASA POWER) \\
      Periodo: & 1 de enero de 2015 al 31 de diciembre de 2024
    \end{tabular}

    \vfill
    \textcolor{brandorange}{\rule{0.92\textwidth}{1.4pt}}\\[0.55cm]
    {\large Puno -- Perú}\\[0.15cm]
    {\large 2026}
  \end{center}
\end{titlepage}

\hypersetup{pageanchor=true}
\pagenumbering{roman}
\tableofcontents
\clearpage
\listoffigures
\clearpage
\listoftables
\clearpage
\pagenumbering{arabic}

% -----------------------------------------------------------------------------
\section{Resumen}

Se desarrolló un dashboard interactivo para caracterizar la frecuencia, intensidad,
estacionalidad y distribución territorial de las heladas en las trece provincias de la
región Puno. La unidad de análisis es la observación diaria por capital provincial. El
dataset integra \NRegistros{} registros del periodo 2015--2024 y \NVariables{} variables
analíticas. La fuente primaria es ERA5-Land / ERA5 del Servicio de Cambio Climático de
Copernicus, consultada mediante Open-Meteo a una resolución aproximada de 9~km; MERRA-2,
obtenida de NASA POWER, se empleó únicamente para validación cruzada.

El proceso incluyó normalización del esquema, conversión de unidades, validación física,
control de duplicados y continuidad temporal, integración de fuentes y construcción de 23
variables derivadas. La tasa de helada meteorológica (temperatura mínima $\leq 0\,^{\circ}\mathrm{C}$)
fue \TasaHelada{}, mientras que la tasa bajo el umbral agronómico ($\leq 3\,^{\circ}\mathrm{C}$)
alcanzó \TasaHeladaAgronomica{}. El riesgo no fue homogéneo: Carabaya acumuló 1,126 días
con helada, frente a ninguno en Sandia y Yunguyo. Julio concentró la mayor incidencia
mensual (\TasaJulio{}). Los contrastes OLS y Mann--Kendall no identificaron una tendencia
significativa en la ventana de diez años. El análisis multivariante mostró que la primera
componente explica \PCAUno{} de la varianza y que \PCANoventa{} componentes alcanzan el
90\,\%. Los resultados se presentan mediante indicadores, filtros, visualizaciones y
explicaciones orientadas a la planificación agraria diferenciada por territorio.

\noindent\textbf{Palabras clave:} helada, riesgo agroclimático, Puno, ERA5-Land, dashboard,
visualización de datos, análisis multivariante.

\subsection*{Abstract}

An interactive dashboard was developed to characterize the frequency, severity,
seasonality, and spatial distribution of frost across the thirteen provinces of Puno,
Peru. The unit of analysis is a daily observation at each provincial capital. The
analytical dataset contains \NRegistros{} records and \NVariables{} variables for
2015--2024. ERA5-Land / ERA5, provided by the Copernicus Climate Change Service and
accessed through Open-Meteo, is the primary source; MERRA-2 from NASA POWER is used only
for cross-validation. Meteorological frost occurred in \TasaHelada{} of records, whereas
the agronomic threshold yielded \TasaHeladaAgronomica{}. Frost exposure varied markedly
among provinces and peaked in July. Neither OLS nor Mann--Kendall detected a statistically
significant monotonic trend over the ten-year window. PCA, clustering, anomaly detection,
and an explanatory Random Forest complement the descriptive views. The resulting system
turns reproducible climate data into indicators and evidence for territorially targeted
agricultural decisions.

\noindent\textbf{Keywords:} frost, agroclimatic risk, Puno, ERA5-Land, dashboard, data
visualization, multivariate analysis.

% -----------------------------------------------------------------------------
\section{Introducción}

La helada ocurre cuando la temperatura del aire próximo a la superficie desciende hasta
un nivel capaz de afectar tejidos vegetales, agua disponible y condiciones de manejo
pecuario. En el altiplano, su efecto depende tanto de la temperatura mínima como de la
estación, la altitud, la humedad, la nubosidad, el viento y la exposición territorial. Por
ello, un único promedio regional puede ocultar contrastes que son decisivos para la
programación de siembras, la focalización del seguro agrario y la preparación de alertas.

Los reanálisis climáticos permiten estudiar el fenómeno de forma continua allí donde la
red de estaciones es escasa. ERA5-Land combina modelación física y asimilación de
observaciones para producir series espacial y temporalmente coherentes (Muñoz-Sabater
et~al., 2021). Su uso no sustituye la medición en superficie, pero ofrece una base
reproducible para comparar territorios bajo una definición común. La visualización
interactiva añade una segunda capa de valor: permite cambiar umbrales, periodos y unidades
territoriales sin perder la trazabilidad del cálculo.

Este proyecto adopta el principio de que cada gráfico debe responder una pregunta. La
vista general resume la magnitud del problema; el análisis descriptivo muestra
distribuciones y evolución; el módulo multidimensional examina relaciones, factores
latentes, grupos y anomalías; la vista geográfica sitúa el riesgo; y la sección final
explicita hallazgos, recomendaciones y límites. La metodología combina estadística
descriptiva, inferencia, aprendizaje automático explicativo y criterios de diseño visual.

% -----------------------------------------------------------------------------
\section{Planteamiento del problema}

Las heladas constituyen uno de los peligros agroclimáticos más relevantes para los
sistemas productivos altoandinos. Sin embargo, los registros disponibles suelen estar
dispersos en servicios especializados y expresados mediante variables y unidades que no
son inmediatamente utilizables por quienes toman decisiones. Esta distancia entre dato y
acción favorece dos simplificaciones: tratar a Puno como un territorio climáticamente
homogéneo y usar exclusivamente el umbral meteorológico de 0~\textdegree C, aun cuando
muchos cultivos presentan daño potencial por debajo de 3~\textdegree C.

La pregunta central es: \emph{¿cómo se distribuyen en el tiempo y en el territorio la
frecuencia y la intensidad de las heladas en las trece provincias de Puno, y qué factores
climáticos se asocian con ellas durante 2015--2024?} El problema exige una solución que
integre adquisición, depuración, análisis y comunicación. Mostrar datos sin interpretación
no es suficiente: el usuario debe poder identificar la ventana crítica, comparar
provincias, evaluar relaciones y descargar el subconjunto que sustenta cada lectura.

La relevancia práctica se observa en la diferencia entre umbrales: se contabilizaron
\NHeladas{} observaciones bajo el umbral meteorológico y \NHeladasAgronomicas{} bajo el
agronómico. Planificar sólo con el primer criterio reduce el fenómeno visible a cerca de
una fracción $1/\FactorAgronomico{}$ del riesgo térmico potencial para el cultivo.

% -----------------------------------------------------------------------------
\section{Objetivos}

\subsection{Objetivo general}

Diseñar, implementar y documentar un dashboard interactivo que transforme datos
agroclimáticos verificables en indicadores, visualizaciones y conclusiones sobre el riesgo
por heladas en la región Puno durante 2015--2024, con énfasis en la toma de decisiones
territorialmente diferenciada.

\subsection{Objetivos específicos}

\begin{enumerate}
  \item Construir un dataset diario, trazable y sin discontinuidades para las trece
        capitales provinciales, documentando fuente, unidades y controles de calidad.
  \item Medir la frecuencia, intensidad, estacionalidad y variación interanual de las
        heladas meteorológicas y agronómicas.
  \item Comparar provincias, pisos ecológicos y zonas agroecológicas para identificar
        concentraciones territoriales del riesgo.
  \item Evaluar asociaciones entre temperatura mínima, altitud, humedad, nubosidad,
        precipitación, radiación y viento mediante métodos lineales y no paramétricos.
  \item Aplicar PCA, K-Means, agrupamiento jerárquico, detección de anomalías y Random
        Forest como técnicas complementarias, declarando su alcance y sus límites.
\end{enumerate}

\begin{table}[H]
\centering
\caption{Preguntas analíticas y método principal de respuesta}
\label{tab:preguntas}
\small
\begin{tabularx}{\textwidth}{>{\bfseries\centering\arraybackslash}p{1.3cm}X>{\raggedright\arraybackslash}p{4.7cm}}
\toprule
Código & Pregunta & Método principal \\
\midrule
P1 & ¿Cuál es la situación general del riesgo por heladas? & KPI, resumen ejecutivo y perfil anual. \\
P2 & ¿Cómo se distribuye la ocurrencia e intensidad entre provincias y meses? & Barras, histograma, caja y calendario térmico. \\
P3 & ¿Existe una evolución o tendencia significativa? & Serie anual, OLS, Mann--Kendall y pendiente de Sen. \\
P4 & ¿Qué provincias concentran el mayor y menor riesgo? & Ranking, mapa interactivo y comparación territorial. \\
P5 & ¿Qué relación existe entre variables ambientales y temperatura mínima? & Pearson, Spearman, dispersión y VIF. \\
P6 & ¿Existen factores latentes, zonas homogéneas y días atípicos? & PCA, K-Means, Ward, Isolation Forest y Random Forest. \\
\bottomrule
\end{tabularx}
\end{table}

% -----------------------------------------------------------------------------
\section{Descripción y procedencia del dataset}

\subsection{Fuentes y cobertura}

La fuente primaria es el reanálisis ERA5-Land / ERA5 del European Centre for
Medium-Range Weather Forecasts y el Copernicus Climate Change Service, consultado mediante
la API histórica de Open-Meteo con el modelo \texttt{era5\_seamless} y agregación en la zona
horaria \texttt{America/Lima}. La fuente secundaria es MERRA-2, distribuida por NASA POWER.
Ambas cubren las mismas fechas y capitales; la segunda no alimenta los indicadores, sino
que permite evaluar si la señal binaria de helada depende de un único producto.

\begin{table}[H]
\centering
\caption{Comparación de las fuentes climáticas utilizadas}
\label{tab:fuentes}
\small
\begin{tabularx}{\textwidth}{>{\bfseries\raggedright\arraybackslash}p{3.2cm}XX}
\toprule
Criterio & Fuente primaria & Fuente secundaria \\
\midrule
Producto & ERA5-Land / ERA5 & MERRA-2 \\
Institución & ECMWF--Copernicus C3S; acceso Open-Meteo & NASA Langley Research Center; NASA POWER \\
Resolución nominal & 0,1\textdegree{}, aproximadamente 9~km & 0,5\textdegree{} $\times$ 0,625\textdegree{}, aproximadamente 55--65~km \\
Granularidad & Diaria por capital provincial & Diaria por capital provincial \\
Función analítica & Fuente de todos los indicadores y modelos & Validación cruzada independiente \\
Periodo & 2015-01-01 a 2024-12-31 & 2015-01-01 a 2024-12-31 \\
\bottomrule
\end{tabularx}
\end{table}

La elección de la malla más fina responde a una prueba previa. Con MERRA-2, Chucuito, El
Collao y Yunguyo recibían series idénticas al caer en una misma celda; lo mismo ocurría con
Azángaro, Huancané y San Antonio de Putina. Además, la celda asignaba 3,668~msnm a
Sandia, cuya capital se ubica alrededor de 2,170--2,300~msnm. Un ranking elaborado sobre
esa malla habría confundido la geometría del producto con diferencias provinciales. ERA5-
Land distingue los trece puntos, con elevaciones de malla entre 2,294 y 4,321~msnm.

\subsection{Estructura y cumplimiento de los requisitos}

La clave natural es el par \texttt{(provincia, fecha)}. Existen \NDias{} días por provincia,
\NProvincias{} provincias y, por multiplicación, \NRegistros{} registros. Las \NVariables{}
columnas incluyen variables temporales, territoriales, térmicas, hídricas, de radiación y
viento, indicadores derivados y campos de control de calidad.

\begin{table}[H]
\centering
\caption{Verificación de los requisitos oficiales del conjunto de datos}
\label{tab:requisitos-datos}
\small
\begin{tabularx}{\textwidth}{X>{\centering\arraybackslash}p{2cm}X}
\toprule
Requisito & Estado & Evidencia \\
\midrule
Al menos 500 registros & Cumple & \NRegistros{} observaciones. \\
Al menos seis variables & Cumple & \NVariables{} variables documentadas. \\
Variables numéricas y categóricas & Cumple & 31 numéricas, 14 categóricas y 6 booleanas; fecha temporal. \\
Dimensión temporal o geográfica & Cumple & Fecha diaria, provincia, latitud, longitud y altitud. \\
Fuente verificable y citada & Cumple & ERA5-Land / ERA5 y MERRA-2, con DOI o portal institucional. \\
Necesidad de limpieza o integración & Cumple & Pipeline de 13 etapas, cuatro conversiones e integración de 47,489 pares. \\
Sin información sensible & Cumple & Sólo variables ambientales agregadas por punto de malla. \\
\bottomrule
\end{tabularx}
\end{table}

% -----------------------------------------------------------------------------
\section{Metodología de limpieza y transformación}

\subsection{Pipeline reproducible}

La ingeniería de datos se implementó como funciones independientes de Streamlit. El mismo
pipeline puede invocarse desde scripts, pruebas o cuadernos, lo que reduce divergencias
entre la aplicación y el informe. La Tabla~\ref{tab:pipeline} resume sus trece etapas.

\begin{longtable}{>{\centering\arraybackslash}p{0.8cm}>{\raggedright\arraybackslash}p{4.2cm}>{\raggedright\arraybackslash}p{9.3cm}}
\caption{Etapas del pipeline de preparación de datos}\label{tab:pipeline}\\
\toprule
\textbf{N.} & \textbf{Etapa} & \textbf{Operación y control} \\
\midrule
\endfirsthead
\toprule
\textbf{N.} & \textbf{Etapa} & \textbf{Operación y control} \\
\midrule
\endhead
1 & Normalización & Renombrado de 19 variables nativas y unificación de campos territoriales. \\
2 & Fecha & Conversión de texto ISO a tipo temporal y orden cronológico. \\
3 & Unidades & Conversión de viento, presión e insolación a unidades analíticas comunes. \\
4 & Centinelas & Sustitución auditable de códigos de ausencia; no se encontraron $-999$. \\
5 & Rangos físicos & Comprobación de límites plausibles en 19 variables; cero valores anulados. \\
6 & Coherencia térmica & Validación de mínima $\leq$ media $\leq$ máxima; cero incoherencias. \\
7 & Precipitación & Marcado, sin borrado, de acumulados superiores a 100~mm/día; un caso. \\
8 & Unicidad & Eliminación de duplicados sobre provincia--fecha; cero duplicados. \\
9 & Calendario & Reindexado a \NDias{} días por provincia; cero fechas ausentes. \\
10 & Imputación & Interpolación intraprovincial sólo cuando fuese necesaria; cero imputaciones. \\
11 & Integración & Unión uno a uno con MERRA-2 y cálculo de concordancia. \\
12 & Derivación & Generación de 23 campos temporales, territoriales y de riesgo. \\
13 & Tipado & Categorías y flotantes de 32 bits; 6,8~MB de memoria analítica. \\
\bottomrule
\end{longtable}

\subsection{Conversión de unidades y control de calidad}

\begin{table}[H]
\centering
\caption{Conversiones de unidad aplicadas antes del análisis}
\label{tab:conversiones}
\small
\begin{tabular}{lll}
\toprule
Variable & Unidad original & Unidad final \\
\midrule
Velocidad del viento & km/h & m/s, factor $0{,}277778$ \\
Racha máxima & km/h & m/s, factor $0{,}277778$ \\
Presión superficial & hPa & kPa, factor $0{,}1$ \\
Duración de insolación & s & h, factor $0{,}000277778$ \\
\bottomrule
\end{tabular}
\end{table}

La radiación solar requirió una revisión específica. El servicio entrega el acumulado en
MJ/m$^2$/día. Interpretarlo como kWh/m$^2$/día y limitarlo al intervalo 0--12 habría
invalidado 46,564 observaciones, que luego podrían haber sido rellenadas por interpolación
sin que el error fuese visible. El rango correcto, 0--45~MJ/m$^2$/día, no anuló ninguna
observación. Este episodio demuestra que la validación debe basarse en la unidad declarada
por el servicio y no sólo en el nombre del parámetro.

Se conservó un registro con precipitación superior a 100~mm el 10 de agosto de 2015 y se
marcó mediante \texttt{precipitacion\_sospechosa}. La decisión evita eliminar un episodio
posiblemente real y permite incluirlo o excluirlo con trazabilidad. El dataset final no
contiene ausentes, duplicados ni fechas faltantes.

\subsection{Definiciones operativas e índice de riesgo}

Se definió helada meteorológica cuando $T_{\min}\leq 0\,^{\circ}\mathrm{C}$ y helada agronómica
cuando $T_{\min}\leq 3\,^{\circ}\mathrm{C}$. La intensidad meteorológica se dividió en ligera
(0 a $-2\,^{\circ}\mathrm{C}$), moderada ($-2$ a $-4\,^{\circ}\mathrm{C}$), severa ($-4$ a
$-6\,^{\circ}\mathrm{C}$) y extrema (menor que $-6\,^{\circ}\mathrm{C}$).

El índice diario de riesgo de helada, en escala 0--100, combina cuatro mecanismos. Todos
se normalizan con límites fijos y se recortan a $[0,1]$:

\begin{align}
D &= \operatorname{clip}\!\left(\frac{3-T_{\min}}{12},0,1\right), &
C &= \operatorname{clip}\!\left(1-\frac{N}{100},0,1\right),\\
S &= \operatorname{clip}\!\left(\frac{100-H}{60},0,1\right), &
V &= \operatorname{clip}\!\left(1-\frac{v}{8},0,1\right),\\
IRH &= 100\left(0{,}45D+0{,}25C+0{,}20S+0{,}10V\right),
\label{eq:riesgo}
\end{align}

donde $N$ es nubosidad porcentual, $H$ humedad relativa y $v$ velocidad del viento en
m/s. El déficit térmico recibe el mayor peso por representar el efecto directo; el cielo
despejado, la sequedad y la calma describen condiciones favorables a la pérdida radiativa.
Los límites fijos hacen comparables dos selecciones filtradas. Los pesos son explícitos,
pero no están calibrados contra pérdidas de rendimiento y por ello el índice no debe
interpretarse como probabilidad de daño.

\subsection{Validación cruzada}

Los \NRegistros{} pares ERA5-Land--MERRA-2 presentaron correlación $r=\FuenteR{}$, sesgo
medio de \FuenteSesgo{} y error absoluto medio de \FuenteMAE{}. La clasificación binaria
de helada coincidió en \FuenteConcordancia{}. La concordancia respalda la señal general;
la correlación moderada es coherente con la diferencia entre mallas de 9 y 55--65~km y
obliga a interpretar con más confianza los patrones relativos que los valores puntuales.

% -----------------------------------------------------------------------------
\section{Diseño del dashboard}

\subsection{Arquitectura por capas}

La aplicación separa configuración, dominio, componentes y vistas. Esta decisión reduce
el acoplamiento y permite probar los cálculos sin iniciar un servidor web.

\begin{table}[H]
\centering
\caption{Responsabilidad de las capas del proyecto}
\label{tab:arquitectura}
\small
\begin{tabularx}{\textwidth}{>{\ttfamily}p{3.0cm}X}
\toprule
Directorio & Responsabilidad \\
\midrule
config/ & Rutas, metadatos, fuentes, umbrales, provincias, preguntas e identidad visual. \\
utils/ & Adquisición, limpieza, filtros, métricas, estadística, modelos y exportación. \\
components/ & Tarjetas, KPI, gráficos, barra lateral, navegación y descargas reutilizables. \\
pages/ & Orquestación de las vistas; consume funciones de dominio sin duplicar cálculos. \\
scripts/ & Flujo reproducible de descarga, procesamiento, resultados y capturas. \\
data/ & Originales, dataset analítico, diccionario y metadatos de trazabilidad. \\
report/ & Fuente LaTeX, resultados calculados y figuras del informe. \\
\bottomrule
\end{tabularx}
\end{table}

\subsection{Diseño visual y accesibilidad}

La interfaz usa una identidad pastel de naranja y azul en superficies amplias. Los trazos,
marcadores y textos emplean versiones más oscuras para conservar contraste. La paleta no
se seleccionó únicamente por preferencia estética: se optimizó en el espacio OKLCH bajo
restricciones de claridad, croma, contraste con el fondo y separación cromática simulada
para protanopía y deuteranopía según el modelo de Machado, Oliveira y Fernandes (2009).

Se aplicaron tres reglas. Primero, una categoría nominal comparada en barras mantiene un
color común y usa posición y longitud como canales principales; el color destaca sólo una
referencia. Segundo, las variables continuas usan escalas secuenciales y las cantidades con
signo, escalas divergentes centradas en cero. Tercero, toda codificación cromática dispone
de un canal de relevo mediante etiquetas, marcadores o tabla descargable. La dispersión y
los mapas limitan a cuatro series cuando cualquier par de marcas puede quedar adyacente.

\subsection{Vistas, filtros e interacción}

La vista general y cinco páginas analíticas conservan un estado compartido. Los ocho
controles incluyen rango temporal, provincias, pisos ecológicos, temporadas, temperatura
mínima, altitud, umbral de helada y exclusión de registros sospechosos. Al modificar un
control se recalculan KPI, gráficos, tablas e interpretaciones. La selección persiste al
navegar, lo que constituye filtrado cruzado real.

\begin{table}[H]
\centering
\caption{Módulos del dashboard y propósito analítico}
\label{tab:modulos}
\small
\begin{tabularx}{\textwidth}{>{\bfseries}p{3.4cm}X}
\toprule
Vista & Contenido y propósito \\
\midrule
Vista general & Ocho KPI, resumen ejecutivo, hallazgos automáticos, riesgo, ranking y ciclo anual. \\
Análisis descriptivo & Distribuciones, comparación provincial, estacionalidad y tendencia P2--P4. \\
Análisis multidimensional & Correlación, PCA, agrupamientos, anomalías y modelo explicativo P5--P6. \\
Análisis geográfico & Mapa interactivo, perfil altitudinal, cuencas, pisos y efecto circunlacustre. \\
Conclusiones & Respuesta explícita a P1--P6, recomendaciones, límites, fuentes y trazabilidad. \\
Metodología y datos & Fuentes, requisitos, bitácora, unidades, diccionario, validación y descargas. \\
\bottomrule
\end{tabularx}
\end{table}

El inventario supera las seis visualizaciones obligatorias: barras para ranking; líneas
para evolución; histograma y caja para distribución térmica; dispersión para relaciones;
mapa de calor estacional; mapa geográfico; scree plot, biplot y proyección 3D para PCA;
dendrograma para Ward; cronología de anomalías; curvas ROC y precisión--recobrado para el
clasificador. Cada gráfico se acompaña de descripción, patrón, hallazgo, decisión y
limitación. Los datos filtrados pueden descargarse en CSV con codificación UTF-8 BOM y en
un libro Excel de varias hojas.

% -----------------------------------------------------------------------------
\section{Resultados e interpretación}

Los resultados de esta sección se recalcularon desde \path{data/datos_limpios.csv} con
semilla 42 y se almacenaron en \path{report/resultados.json}. Las figuras se generan con
\path{report/generar_figuras_informe.py}; no son transcripciones manuales.

\subsection{P1: situación general del riesgo}

De \NRegistros{} observaciones, \NHeladas{} cumplieron el umbral meteorológico, equivalente
a \TasaHelada{}. El umbral agronómico identificó \NHeladasAgronomicas{} observaciones
(\TasaHeladaAgronomica{}), \FactorAgronomico{} veces la incidencia meteorológica. La
temperatura mínima media fue \TminMedia{} y la mínima absoluta, \TminAbsoluta{} en
Carabaya el 23 de julio de 2018. El índice compuesto promedió \RiesgoMedio{} puntos.

\begin{table}[H]
\centering
\caption{Distribución de la intensidad de helada meteorológica}
\label{tab:intensidad}
\begin{tabular}{lrr}
\toprule
Clase & Observaciones & Proporción del total \\
\midrule
Sin helada & 39,869 & 83,95\,\% \\
Ligera & 5,180 & 10,91\,\% \\
Moderada & 2,090 & 4,40\,\% \\
Severa & 320 & 0,67\,\% \\
Extrema & 30 & 0,06\,\% \\
\bottomrule
\end{tabular}
\end{table}

La mayoría de los eventos fue ligera o moderada; aun así, 350 observaciones alcanzaron
intensidad severa o extrema. La decisión operativa no debe basarse sólo en el promedio:
los extremos y el umbral agronómico cambian el volumen de días que requieren prevención.

\subsection{P2: distribución provincial y ciclo anual}

\projectfigure{figuras/01_ranking_provincial.png}{Ranking provincial de días con helada meteorológica. La longitud de la barra codifica el conteo acumulado bajo un periodo común para todas las provincias.}{fig:ranking}

Carabaya encabezó el ranking con 1,126 días (30,82\,\% de sus observaciones), seguida por
San Antonio de Putina con 976 y El Collao con 962. Yunguyo y Sandia no presentaron días
bajo 0~\textdegree C en el punto de malla analizado. Esta ausencia no implica que toda la
provincia esté libre de heladas: cada territorio se representa mediante una sola capital.

\begin{table}[H]
\centering
\caption{Indicadores térmicos por provincia, 2015--2024}
\label{tab:provincias}
\scriptsize
\begin{tabular}{lrrrr}
\toprule
Provincia & Altitud (m) & Días & Tasa (\%) & $T_{\min}$ media (\textdegree C) \\
\midrule
Carabaya & 4,321 & 1,126 & 30,82 & 0,96 \\
San Antonio de Putina & 3,867 & 976 & 26,72 & 2,24 \\
El Collao & 3,835 & 962 & 26,33 & 2,85 \\
San Román & 3,828 & 930 & 25,46 & 2,53 \\
Azángaro & 3,861 & 830 & 22,72 & 2,45 \\
Lampa & 3,868 & 799 & 21,87 & 2,72 \\
Puno & 3,819 & 683 & 18,70 & 3,41 \\
Melgar & 3,912 & 460 & 12,59 & 2,82 \\
Chucuito & 3,872 & 406 & 11,11 & 4,10 \\
Huancané & 3,838 & 365 & 9,99 & 3,93 \\
Moho & 4,043 & 83 & 2,27 & 4,05 \\
Sandia & 2,294 & 0 & 0,00 & 10,10 \\
Yunguyo & 3,825 & 0 & 0,00 & 7,15 \\
\bottomrule
\end{tabular}
\end{table}

\projectfigure{figuras/02_ciclo_anual.png}{Ciclo anual promedio. Los paneles separados evitan un doble eje y muestran que el máximo de heladas coincide con la estación de menor precipitación.}{fig:ciclo}

La incidencia fue menor que 1\,\% entre enero y marzo y alcanzó 47,44\,\% en junio y
\TasaJulio{} en julio. En esos meses la precipitación diaria media descendió a 0,79 y
0,63~mm, respectivamente. El patrón es consistente con heladas radiativas bajo cielos más
despejados y aire seco; la visualización identifica mayo--agosto como la ventana prioritaria
para preparación, sin afirmar que precipitación y helada mantengan una relación causal
aislada.

\projectfigure{figuras/07_calendario_heladas.png}{Calendario provincia--mes de incidencia de heladas. Las celdas muestran el porcentaje de observaciones bajo 0~\textdegree C.}{fig:calendario}

\subsection{P3: evolución y tendencia}

\projectfigure{figuras/03_tendencia_anual.png}{Promedio anual de días con helada por provincia y ajuste lineal OLS.}{fig:tendencia}

La tasa anual osciló entre 11,14\,\% en 2024 y 18,23\,\% en 2015. OLS estimó
\OLSPendiente{} días por provincia y año, pero el intervalo de confianza al 95\,\% incluyó
cero y $p=\OLSp{}$ ($R^2=\OLSRdos{}$). Mann--Kendall produjo
$\tau=\MKTau{}$, pendiente de Sen \SenPendiente{} y $p=\MKp{}$. Ambos métodos coinciden:
no existe evidencia suficiente de tendencia monotónica al nivel 0,05.

\begin{table}[H]
\centering
\caption{Contrastes de tendencia sobre los días anuales de helada por provincia}
\label{tab:tendencia}
\begin{tabular}{lrrrrl}
\toprule
Método & Pendiente & Estadístico & $p$ & $n$ & Lectura \\
\midrule
OLS & -1,37 días/año & $R^2=0,260$ & 0,132 & 10 & No significativa \\
Mann--Kendall/Sen & -1,00 días/año & $\tau=-0,422$ & 0,107 & 10 & Sin tendencia \\
\bottomrule
\end{tabular}
\end{table}

La caída observada en 2024 no se interpreta como cambio climático. Diez años son
insuficientes para establecer una normal climática y la variabilidad interanual puede
dominar el ajuste. El resultado sí permite afirmar que el periodo estudiado no ofrece
evidencia estadística de una reducción sostenida.

\subsection{P4: concentración territorial del riesgo}

El contraste más informativo aparece entre Yunguyo y San Román. Sus elevaciones de malla
son casi idénticas (3,825 y 3,828~msnm), pero acumularon 0 y 930 días con helada. En el
conjunto de trece provincias, la relación entre altitud y tasa fue positiva
($r=\AltitudR{}$), aunque no significativa al 5\,\% ($p=\AltitudP{}$). En el nivel diario,
la correlación entre altitud y temperatura mínima fue $r=-0,551$.

Las observaciones de la zona ``Altiplano circunlacustre'' presentaron una tasa de
\LagoTasa{}, frente a \RestoTasa{} en el resto. La diferencia es consistente con el efecto
termorregulador del lago Titicaca, cuya inercia térmica atenúa el enfriamiento nocturno.
No obstante, el diseño observacional y la representación por un punto no permiten aislar
causalmente el efecto del lago frente a topografía, cobertura y circulación local. La
decisión prudente es evitar zonificaciones basadas sólo en altitud y combinar la posición
circunlacustre con evidencia local.

\subsection{P5: relación entre variables agroclimáticas}

\projectfigure[0.82\textwidth]{figuras/04_correlaciones.png}{Correlaciones de Pearson entre variables seleccionadas. La escala divergente separa asociaciones negativas y positivas.}{fig:correlaciones}

La temperatura mínima mostró las asociaciones lineales más altas con temperatura del
suelo ($r=0,658$), nubosidad ($r=0,561$), altitud ($r=-0,551$) y humedad relativa
($r=0,536$). La precipitación alcanzó $r=0,437$ y la radiación solar $r=-0,105$. Spearman
confirmó el orden general, con $\rho=0,679$ para temperatura del suelo y $\rho=0,557$ para
precipitación. Dado el tamaño muestral, magnitudes pequeñas pueden resultar
estadísticamente detectables; por eso la interpretación prioriza tamaño y mecanismo, no
sólo el valor $p$.

La asociación positiva entre humedad, nubosidad y mínima diaria es físicamente coherente
con la reducción de la pérdida radiativa nocturna. Sin embargo, una correlación bivariada
no separa estación, provincia ni colinealidad. El dashboard ofrece Pearson y Spearman,
valores $p$ y VIF, y usa la dispersión para reconocer no linealidad y grupos territoriales.

\subsection{P6: factores latentes, agrupamientos, anomalías y explicación}

\projectfigure{figuras/05_pca_varianza.png}{Varianza individual y acumulada del PCA aplicado a doce variables estandarizadas.}{fig:pca}

El PCA se ejecutó después de estandarizar las doce variables; de otro modo la altitud,
medida en miles de metros, dominaría por escala. CP1 explicó \PCAUno{} y estuvo definida
principalmente por humedad relativa y punto de rocío con cargas positivas, frente a
oscilación térmica negativa: un eje humedad--amplitud. CP2 explicó 20,46\,\% y concentró
temperatura del suelo, evapotranspiración, radiación y temperatura mínima: un eje térmico-
energético. Las dos primeras sumaron \PCADosAcum{}; fueron necesarias \PCANoventa{}
componentes para superar 90\,\%, por lo que reducir todo el sistema a dos dimensiones
facilita la visualización pero no conserva toda la información.

K-Means seleccionó $k=\KOptimo{}$ al maximizar la silueta entre dos y cuatro grupos. La
silueta fue \Silueta{}, clasificada como estructura débil. Los perfiles separaron un
régimen seco con helada frecuente (34,4\,\% y riesgo medio 47,44), uno húmedo con baja
incidencia (3,2\,\%) y el régimen cálido de Sandia (0\,\%). Los grupos son una síntesis
exploratoria, no fronteras agroclimáticas definitivas.

Isolation Forest marcó \NAnomalias{} observaciones, aproximadamente 1\,\% del total. Las
más extremas combinaron altas temperaturas, lluvia intensa, humedad cercana a saturación,
nubosidad y baja radiación, especialmente en Sandia. Esta lectura ilustra la ventaja del
criterio multivariante: una observación puede ser atípica por combinación aunque ninguna
variable aislada supere un umbral fijo.

El Random Forest explicativo, entrenado con partición cronológica, alcanzó AUC ROC
\RFAUC{}, precisión media \RFAP{}, exactitud \RFExactitud{}, sensibilidad
\RFSensibilidad{} y precisión positiva \RFPrecision{}. Se excluyeron temperatura mínima,
déficit térmico e índice de riesgo porque definen directa o algebraicamente el objetivo.
La importancia por permutación situó primero a oscilación térmica, punto de rocío y
temperatura del suelo. El desempeño no convierte el modelo en pronóstico operativo: no se
ha validado prospectivamente ni contra estaciones, y su propósito es ordenar variables
asociadas bajo el periodo disponible.

\projectfigure[0.72\textwidth]{figuras/06_validacion_fuentes.png}{Comparación de temperatura mínima entre ERA5-Land y MERRA-2. La línea discontinua representa concordancia perfecta.}{fig:validacion}

% -----------------------------------------------------------------------------
\section{Capturas del dashboard}

Las capturas documentan la traducción de los resultados a una interfaz navegable. Se
generan automáticamente con un navegador sin interfaz, resolución 1680$\times$1050 y
factor de escala 2; por tanto, constituyen evidencia reproducible y no una composición
manual. Si una imagen no está disponible durante una compilación intermedia, el documento
mantiene el espacio identificado sin interrumpir el resto del informe.

\projectfigure{../assets/capturas/01_vista_general.png}{Vista general: identificación, filtros, KPI y síntesis del periodo activo.}{fig:captura-general}

La primera vista debe permitir responder P1 sin navegar: explicita denominador, umbral,
periodo y provincia crítica. Las variaciones se colorean por efecto favorable o adverso,
no sólo por signo numérico.

\projectfigure{../assets/capturas/02_analisis_descriptivo.png}{Análisis descriptivo: distribuciones, comparación territorial y evolución temporal.}{fig:captura-descriptivo}

El módulo descriptivo vincula cada gráfico con una tabla y una interpretación. Las
distribuciones muestran la forma y los extremos; el ranking conserva el mismo periodo
para evitar comparaciones con denominadores distintos.

\projectfigure{../assets/capturas/03_analisis_multidimensional.png}{Análisis multidimensional: correlaciones, PCA, grupos, anomalías y modelo explicativo.}{fig:captura-multidimensional}

La vista avanzada separa exploración, diagnóstico y desempeño. Los controles limitan el
número de series y mantienen visibles la estandarización, la semilla y los criterios de
selección de componentes o grupos.

\projectfigure{../assets/capturas/04_analisis_geografico.png}{Análisis geográfico: localización provincial, perfil altitudinal y calendario territorial.}{fig:captura-geografico}

El mapa permite ubicar el riesgo, pero el ranking y las etiquetas aportan la comparación
cuantitativa que el área y el color por sí solos no garantizan. La vista declara que cada
provincia está representada por su capital.

\projectfigure{../assets/capturas/05_conclusiones.png}{Conclusiones: respuestas a P1--P6, recomendaciones y limitaciones.}{fig:captura-conclusiones}

Esta página evita que el usuario deba inferir conclusiones sólo a partir de gráficos. Cada
respuesta cita la magnitud relevante, enlaza el método y separa resultado, decisión y
limitación.

\projectfigure{../assets/capturas/06_metodologia_datos.png}{Metodología y datos: fuentes, bitácora ETL, diccionario y descargas.}{fig:captura-metodologia}

La trazabilidad muestra el origen de los datos, las conversiones, los controles y la
validación. Los archivos original, limpio, diccionario y metadatos se mantienen separados
para que el resultado pueda auditarse.

% -----------------------------------------------------------------------------
\section{Conclusiones}

\begin{enumerate}
  \item La incidencia depende del umbral: \TasaHelada{} de las observaciones cumplió la
        definición meteorológica y \TasaHeladaAgronomica{} la agronómica. Usar sólo
        0~\textdegree C subrepresenta el conjunto de días térmicamente relevantes para el
        cultivo por un factor de \FactorAgronomico{}.
  \item El riesgo no es homogéneo. Carabaya acumuló 1,126 días con helada, mientras Sandia
        y Yunguyo no registraron ninguno en el punto analizado. Una medida regional uniforme
        oculta diferencias de más de treinta puntos porcentuales.
  \item La ventana crítica se concentra entre mayo y agosto; julio alcanzó
        \TasaJulio{}. Este calendario permite anticipar la asignación de medidas de
        protección, alertas y asistencia técnica.
  \item La altitud es relevante pero insuficiente. Yunguyo y San Román tienen elevaciones
        de malla casi idénticas y resultados opuestos. La menor tasa circunlacustre es
        consistente con la inercia térmica del Titicaca y exige incorporar ubicación y
        microclima a la zonificación.
  \item No se detectó una tendencia monotónica significativa durante 2015--2024: OLS produjo
        $p=\OLSp{}$ y Mann--Kendall $p=\MKp{}$. El hallazgo correcto es ausencia de evidencia
        en esta ventana, no demostración de estabilidad climática de largo plazo.
  \item La estructura multivariante no se reduce sin pérdida a dos ejes. CP1 y CP2 explican
        \PCADosAcum{}, pero se requieren \PCANoventa{} componentes para 90\,\%. La
        proyección 2D es útil para explorar, no para sustituir las variables originales.
  \item K-Means obtuvo tres perfiles con silueta \Silueta{}. La separación débil obliga a
        usar los grupos como tipología exploratoria y no como límites administrativos o
        ecológicos rígidos.
  \item La validación cruzada logró \FuenteConcordancia{} en la clasificación de helada. La
        señal general es reproducible entre productos, aunque las diferencias de resolución
        justifican cautela en valores absolutos.
  \item El producto final es reproducible: conserva original, limpio, diccionario,
        metadatos, semilla y código de generación. El dashboard no sólo muestra figuras;
        conecta filtros, resultados, interpretación y descarga del subconjunto analizado.
\end{enumerate}

% -----------------------------------------------------------------------------
\section{Recomendaciones}

\begin{description}[leftmargin=4.0cm,style=nextline,itemsep=0.65em]
  \item[Gobierno Regional de Puno.] Priorizar seguro agrario, cobertizos, riego preventivo
        y asistencia en Carabaya, San Antonio de Putina, El Collao y San Román, sin asumir
        que una misma intervención produce igual beneficio en toda la región.
  \item[Dirección Regional Agraria.] Elaborar calendarios de siembra y protocolos de
        alerta por provincia. Intensificar la vigilancia entre mayo y agosto y complementar
        temperatura prevista con humedad, nubosidad y oscilación térmica.
  \item[Productores y organizaciones.] Usar el umbral agronómico de 3~\textdegree C como
        aviso temprano para cultivos sensibles; combinarlo con conocimiento local, estado
        fenológico y disponibilidad de riego o cobertura.
  \item[SENAMHI y academia.] Densificar estaciones en fondos de valle, laderas y zonas con
        mayor discrepancia entre reanálisis. Comparar prospectivamente el índice con
        observaciones de superficie y eventos reportados.
  \item[Equipo de desarrollo.] Mantener pruebas automáticas sobre filtros, KPI, agregaciones
        y modelos; versionar el esquema del dataset y recalcular \texttt{resultados.json}
        antes de cada publicación.
  \item[Investigación futura.] Integrar rendimiento y pérdidas económicas, cobertura de
        suelo, pendiente, distancia al lago y pronósticos de corto plazo. Calibrar los pesos
        del índice contra daño observado y evaluar validación espacial fuera de las capitales.
\end{description}

% -----------------------------------------------------------------------------
\section{Limitaciones}

\begin{enumerate}
  \item La malla primaria de aproximadamente 9~km suaviza extremos locales y no resuelve
        fondos de valle ni laderas donde la inversión térmica puede intensificarse.
  \item Los datos son reanálisis, es decir, combinación de modelo y observaciones
        asimiladas; no equivalen a una estación meteorológica de superficie.
  \item Cada provincia se representa mediante la coordenada de su capital. La variabilidad
        interna de Carabaya, Sandia y otros territorios topográficamente complejos queda
        submuestreada.
  \item Diez años permiten estudiar variabilidad interanual, pero no establecer normales
        climáticas de treinta años ni atribuir tendencias al cambio climático.
  \item MERRA-2 posee una malla más gruesa. Su correlación moderada con ERA5-Land refleja
        escalas distintas y limita la comparación de valores diarios puntuales.
  \item Un registro de precipitación superior a 100~mm fue conservado y marcado; puede
        influir en análisis sensibles a extremos, aunque es excluible mediante filtro.
  \item No se dispone de rendimiento agrícola, especie, variedad ni fase fenológica; por
        ello no se cuantifica daño físico o económico.
  \item El índice de riesgo usa pesos expertos y límites fijos no calibrados contra pérdidas
        observadas. Es una herramienta comparativa, no una probabilidad validada de daño.
  \item El Random Forest es explicativo y fue evaluado dentro del mismo periodo. La alta
        AUC no demuestra capacidad de pronóstico futuro ni transferencia a otros lugares.
  \item Las asociaciones y la comparación circunlacustre son observacionales. No permiten
        inferir causalidad aislada ni separar todos los mecanismos topográficos y atmosféricos.
\end{enumerate}

% -----------------------------------------------------------------------------
\section{Referencias}

\begin{enumerate}[label={},leftmargin=1.2em,itemindent=-1.2em,itemsep=0.65em]
  \item Breiman, L. (2001). Random forests. \emph{Machine Learning, 45}(1), 5--32.
        \url{https://doi.org/10.1023/A:1010933404324}
  \item Cleveland, W. S., y McGill, R. (1984). Graphical perception: Theory,
        experimentation, and application to the development of graphical methods.
        \emph{Journal of the American Statistical Association, 79}(387), 531--554.
  \item Cohen, J. (1988). \emph{Statistical power analysis for the behavioral sciences}
        (2.a ed.). Lawrence Erlbaum Associates.
  \item Gelaro, R., et al. (2017). The Modern-Era Retrospective Analysis for Research and
        Applications, Version 2 (MERRA-2). \emph{Journal of Climate, 30}(14), 5419--5454.
        \url{https://doi.org/10.1175/JCLI-D-16-0758.1}
  \item Hersbach, H., et al. (2020). The ERA5 global reanalysis. \emph{Quarterly Journal of
        the Royal Meteorological Society, 146}(730), 1999--2049.
        \url{https://doi.org/10.1002/qj.3803}
  \item Iglewicz, B., y Hoaglin, D. C. (1993). \emph{How to detect and handle outliers}.
        ASQC Quality Press.
  \item Jolliffe, I. T., y Cadima, J. (2016). Principal component analysis: A review and
        recent developments. \emph{Philosophical Transactions of the Royal Society A,
        374}(2065). \url{https://doi.org/10.1098/rsta.2015.0202}
  \item Kendall, M. G. (1975). \emph{Rank correlation methods} (4.a ed.). Charles Griffin.
  \item Liu, F. T., Ting, K. M., y Zhou, Z.-H. (2008). Isolation forest. En \emph{2008
        Eighth IEEE International Conference on Data Mining} (pp. 413--422).
        \url{https://doi.org/10.1109/ICDM.2008.17}
  \item Machado, G. M., Oliveira, M. M., y Fernandes, L. A. F. (2009). A
        physiologically-based model for simulation of color vision deficiency.
        \emph{IEEE Transactions on Visualization and Computer Graphics, 15}(6), 1291--1298.
  \item Mann, H. B. (1945). Nonparametric tests against trend. \emph{Econometrica, 13}(3),
        245--259.
  \item Muñoz-Sabater, J., et al. (2021). ERA5-Land: A state-of-the-art global reanalysis
        dataset for land applications. \emph{Earth System Science Data, 13}(9), 4349--4383.
        \url{https://doi.org/10.5194/essd-13-4349-2021}
  \item Munzner, T. (2014). \emph{Visualization analysis and design}. CRC Press.
  \item NASA Langley Research Center. (2025). \emph{NASA Prediction Of Worldwide Energy
        Resources (POWER) Project: Daily agroclimatology data} (Version 2.9.6)
        [Conjunto de datos]. \url{https://power.larc.nasa.gov/}
  \item Pulgar Vidal, J. (1981). \emph{Geografía del Perú: Las ocho regiones naturales}
        (8.a ed.). Editorial Universo.
  \item Rousseeuw, P. J. (1987). Silhouettes: A graphical aid to the interpretation and
        validation of cluster analysis. \emph{Journal of Computational and Applied
        Mathematics, 20}, 53--65.
  \item Sen, P. K. (1968). Estimates of the regression coefficient based on Kendall's tau.
        \emph{Journal of the American Statistical Association, 63}(324), 1379--1389.
  \item Servicio Nacional de Meteorología e Hidrología del Perú. (2024). \emph{Monitoreo y
        pronóstico de heladas}. \url{https://www.senamhi.gob.pe/}
  \item Tukey, J. W. (1977). \emph{Exploratory data analysis}. Addison-Wesley.
  \item Universidad Nacional del Altiplano. (2026). \emph{Proyecto final: Dashboard
        interactivo de visualización de datos} [Consigna del curso SIS328].
\end{enumerate}

% -----------------------------------------------------------------------------
\clearpage
\appendix
\section{Enlaces y reproducibilidad}

\begin{description}[leftmargin=4.1cm,style=nextline]
  \item[Repositorio.] \url{https://github.com/jhonaracayo/dashboard-heladas-puno}
  \item[Dashboard publicado.] \url{https://dashboard-heladas-puno.streamlit.app}
  \item[Fuente primaria.] \url{https://archive-api.open-meteo.com/v1/archive}
  \item[Fuente secundaria.] \url{https://power.larc.nasa.gov/api/temporal/daily/point}
\end{description}

Los comandos mínimos de reproducción, ejecutados desde la raíz del repositorio, son:

\begin{verbatim}
python scripts/01_descargar_datos.py
python scripts/02_procesar_datos.py
python report/generar_figuras_informe.py
streamlit run app.py
\end{verbatim}

La descarga no debe repetirse para reproducir el informe mientras se conserven los
originales. La semilla global es 42. El archivo \texttt{requirements.txt} fija las versiones
mínimas verificadas; \texttt{metadata.json} registra cobertura, estructura y cada etapa del
pipeline.

\section{Bitácora íntegra del pipeline}

\begin{enumerate}[leftmargin=*,itemsep=0.35em]
\item Inicio del pipeline: 47,489 filas \ensuremath{\times} 32 columnas de entrada.
\item Normalización del esquema: 19 variables nativas del servicio se renombraron a nombres semánticos en español y se unificó la nomenclatura geográfica (latitud, longitud, altitud).
\item La fecha llegaba como cadena de texto ISO; se convirtió a tipo temporal para habilitar filtros por rango, agregaciones cronológicas y cálculo de tendencias.
\item Homogeneización de unidades en 4 variables: velocidad\_\allowbreak{}viento km/h \ensuremath{\rightarrow} m/s (\ensuremath{\times}0.277778); racha\_\allowbreak{}viento\_\allowbreak{}maxima km/h \ensuremath{\rightarrow} m/s (\ensuremath{\times}0.277778); presion\_\allowbreak{}superficie hPa \ensuremath{\rightarrow} kPa (\ensuremath{\times}0.1); horas\_\allowbreak{}sol s \ensuremath{\rightarrow} h (\ensuremath{\times}0.000277778).
\item No se hallaron valores centinela (-999.0); el servicio codifica la ausencia como nulo JSON, ya interpretado como ausente en la carga.
\item Validación de rangos físicos sobre 19 variables: 0 observaciones fuera de los límites admisibles fueron anuladas.
\item Verificación de coherencia térmica (mínima \ensuremath{\leq} media \ensuremath{\leq} máxima): 0 filas incoherentes; en ellas la media se recalculó como punto medio de los extremos observados.
\item Control de calidad de la precipitación: 1 acumulados diarios superiores a 100 mm se marcaron como sospechosos y se conservaron sin modificar. El episodio dominante es 2015-08-10, con 1 provincia(s) afectada(s) el mismo día; el dato permanece disponible y trazable para el análisis de anomalías.
\item Control de unicidad sobre la clave natural (provincia, fecha): 0 registros duplicados eliminados.
\item Reindexado al calendario diario completo (3,653 días \ensuremath{\times} 13 provincias = 47,489 filas esperadas): 0 días ausentes materializados.
\item Imputación por provincia (interpolación lineal bidireccional en variables continuas; relleno con cero en las de precipitación): 0 valores completados; 0 ausentes residuales.
\item Integración de la fuente secundaria (MERRA-2 / NASA POWER) por la clave (provincia, fecha): 47,489 filas emparejadas. Correlación de la temperatura mínima entre fuentes r = 0.611; sesgo medio +1.29 \textdegree{}C; concordancia en la clasificación binaria de helada 83.1 \%.
\item Transformación: se generaron 23 variables derivadas (descomposición temporal, caracterización territorial por piso ecológico, indicadores de helada e índice compuesto de riesgo agroclimático).
\item Optimización de tipos (categóricos y punto flotante de 32 bits): el dataset analítico ocupa 6.8 MB en memoria.
\item Pipeline finalizado: 47,489 filas \ensuremath{\times} 52 columnas en el dataset analítico.
\end{enumerate}

\section{Diccionario completo de variables}

El diccionario siguiente se genera desde \texttt{data/diccionario\_datos.csv}. Para cada
campo se indica etiqueta, unidad, procedencia y definición operativa; por ello funciona
como contrato semántico entre adquisición, análisis y visualización.

\small
\begin{longtable}{>{\raggedright\arraybackslash}p{2.8cm}>{\raggedright\arraybackslash}p{2.9cm}>{\centering\arraybackslash}p{1.2cm}>{\centering\arraybackslash}p{1.5cm}>{\raggedright\arraybackslash}p{5.5cm}}
\caption{Diccionario completo del dataset analítico}\label{tab:diccionario-completo}\\
\toprule
\textbf{Variable} & \textbf{Etiqueta} & \textbf{Unidad} & \textbf{Origen} & \textbf{Descripción} \\
\midrule
\endfirsthead
\multicolumn{5}{l}{\small\itshape Continuación de la tabla~\ref{tab:diccionario-completo}}\\
\toprule
\textbf{Variable} & \textbf{Etiqueta} & \textbf{Unidad} & \textbf{Origen} & \textbf{Descripción} \\
\midrule
\endhead
\midrule \multicolumn{5}{r}{\small Continúa en la página siguiente}\\
\endfoot
\bottomrule
\endlastfoot
fecha & Fecha & — & original & Día de observación (resolución diaria). \\
anio & Año & — & derivada & Año calendario de la observación. \\
mes & Mes & — & derivada & Número de mes (1–12). \\
mes\_\allowbreak{}nombre & Nombre del mes & — & derivada & Mes en texto, ordenado cronológicamente. \\
mes\_\allowbreak{}abrev & Mes abrev & — & derivada & Variable auxiliar del modelo de datos. \\
trimestre & Trimestre & — & derivada & Trimestre calendario (T1–T4). \\
dia\_\allowbreak{}anio & Día del año & — & derivada & Día juliano (1–366). \\
semana\_\allowbreak{}anio & Semana anio & — & derivada & Variable auxiliar del modelo de datos. \\
anio\_\allowbreak{}mes & Anio mes & — & derivada & Variable auxiliar del modelo de datos. \\
temporada & Temporada & — & derivada & Lluviosa, Transición o Seca. \\
campania\_\allowbreak{}agricola & Campaña agrícola & — & derivada & Campaña MIDAGRI de agosto a julio. \\
provincia & Provincia & — & original & Provincia de la región Puno. \\
capital & Capital provincial & — & original & Capital que define el punto de consulta. \\
latitud & Latitud & \textdegree{} & original & Latitud decimal del punto de malla. \\
longitud & Longitud & \textdegree{} & original & Longitud decimal del punto de malla. \\
altitud & Altitud & msnm & original & Elevación del punto reportada por la API. \\
cuenca & Vertiente hidrográfica & — & derivada & Titicaca (endorreica) o Amazonas. \\
zona\_\allowbreak{}agroecologica & Zona agroecológica & — & derivada & Caracterización agroecológica provincial. \\
piso\_\allowbreak{}ecologico & Piso ecológico & — & derivada & Piso altitudinal según Pulgar Vidal (1981). \\
temperatura\_\allowbreak{}minima & Temperatura mínima & \textdegree{}C & original & Temperatura mínima diaria del aire a 2 m; variable central del estudio. \\
temperatura\_\allowbreak{}media & Temperatura media & \textdegree{}C & original & Temperatura del aire a 2 m, promedio diario. \\
temperatura\_\allowbreak{}maxima & Temperatura máxima & \textdegree{}C & original & Temperatura máxima diaria del aire a 2 m. \\
oscilacion\_\allowbreak{}termica & Oscilación térmica & \textdegree{}C & derivada & Diferencia entre temperatura máxima y mínima. \\
temperatura\_\allowbreak{}suelo & Temperatura del suelo & \textdegree{}C & original & Temperatura media del suelo entre 0 y 7 cm de profundidad. \\
punto\_\allowbreak{}rocio & Punto de rocío & \textdegree{}C & original & Temperatura de rocío media a 2 m. \\
humedad\_\allowbreak{}relativa & Humedad relativa & \% & original & Humedad relativa media diaria a 2 m. \\
deficit\_\allowbreak{}presion\_\allowbreak{}vapor & Déficit de presión de vapor & kPa & derivada & VPD calculado con la ecuación de Tetens. \\
deficit\_\allowbreak{}presion\_\allowbreak{}vapor\_\allowbreak{}maximo & Déficit de presión de vapor máximo & kPa & original & Valor máximo diario del déficit de presión de vapor. \\
precipitacion & Precipitación & mm & original & Precipitación total diaria (lluvia y equivalente en agua de nieve). \\
lluvia & Lluvia & mm & original & Fracción líquida de la precipitación diaria. \\
nevada & Nevada & cm & original & Espesor de nieve acumulada en el día. \\
horas\_\allowbreak{}precipitacion & Horas con precipitación & h & original & Número de horas del día con precipitación registrada. \\
humedad\_\allowbreak{}suelo & Humedad del suelo & m\textsuperscript{3}/m\textsuperscript{3} & original & Contenido volumétrico de agua del suelo entre 0 y 7 cm. \\
evapotranspiracion & Evapotranspiración de referencia & mm & original & ET\textsubscript{0} diaria calculada con el método FAO-56 Penman-Monteith. \\
radiacion\_\allowbreak{}solar & Radiación solar & MJ/m\textsuperscript{2} & original & Irradiancia global horizontal acumulada en el día. \\
horas\_\allowbreak{}sol & Horas de sol & h & original & Duración de la insolación directa. \\
nubosidad & Nubosidad & \% & original & Cobertura nubosa media diaria; determinante de la helada de radiación. \\
velocidad\_\allowbreak{}viento & Velocidad del viento & m/s & original & Velocidad media del viento a 10 m. \\
racha\_\allowbreak{}viento\_\allowbreak{}maxima & Racha máxima de viento & m/s & original & Velocidad máxima del viento a 10 m. \\
presion\_\allowbreak{}superficie & Presión superficial & kPa & original & Presión atmosférica media en superficie. \\
helada & Helada meteorológica & — & derivada & Verdadero si la temperatura mínima \ensuremath{\leq} 0 \textdegree{}C. \\
helada\_\allowbreak{}agronomica & Helada agronómica & — & derivada & Verdadero si la temperatura mínima \ensuremath{\leq} 3 \textdegree{}C. \\
intensidad\_\allowbreak{}helada & Intensidad de helada & — & derivada & Clasificación ordinal de severidad del evento. \\
deficit\_\allowbreak{}termico & Déficit térmico & \textdegree{}C & derivada & Grados por debajo del umbral agronómico (0 si no aplica). \\
indice\_\allowbreak{}riesgo\_\allowbreak{}helada & Índice de riesgo de helada & 0–100 & derivada & Índice compuesto de riesgo agroclimático diario. \\
nivel\_\allowbreak{}riesgo & Nivel de riesgo & — & derivada & Clasificación ordinal del índice de riesgo en cinco niveles. \\
dia\_\allowbreak{}lluvia & Día con lluvia & — & derivada & Verdadero si la precipitación \ensuremath{\geq} 1 mm. \\
dia\_\allowbreak{}nevada & Dia nevada & — & derivada & Variable auxiliar del modelo de datos. \\
precipitacion\_\allowbreak{}sospechosa & Precipitación implausible & — & derivada & Marca de control de calidad: acumulado diario > 100 mm, atribuible al reanálisis. \\
temperatura\_\allowbreak{}minima\_\allowbreak{}merra2 & Temperatura mínima (MERRA-2) & \textdegree{}C & original & Temperatura mínima de la fuente secundaria, para validación cruzada. \\
helada\_\allowbreak{}merra2 & Helada según MERRA-2 & — & derivada & Ocurrencia de helada según la fuente secundaria independiente. \\
fuente & Fuente & — & derivada & Variable auxiliar del modelo de datos. \\
\end{longtable}
\normalsize

\section{Aportes del equipo}

\begin{table}[H]
\centering
\caption{Responsabilidades y contribuciones de los integrantes}
\label{tab:aportes}
\small
\begin{tabularx}{\textwidth}{>{\bfseries}p{4.0cm}p{3.3cm}X}
\toprule
Integrante & Rol principal & Aportes \\
\midrule
Aracayo Mamani, Jhon Marco & Arquitectura de software y dashboard & Estructura por capas, integración de Streamlit, navegación, estado de filtros y componentes reutilizables. \\
Canaza Paucara, Juan Diego & Ingeniería de datos y análisis multivariante & Adquisición, pipeline de limpieza, integración MERRA-2, PCA, agrupamiento, anomalías y reproducibilidad. \\
Luque Pacheco, Angie Tatiana & Visualización, estadística y redacción & Diseño visual accesible, selección y justificación de gráficos, análisis de tendencia, interpretación e informe técnico. \\
\bottomrule
\end{tabularx}
\end{table}

La participación se distribuye por especialidad, pero los tres integrantes comparten la
responsabilidad de explicar el problema, la preparación de datos, el código, las
visualizaciones y las conclusiones durante la sustentación.

\end{document}
